\newpage
\chapter{CONCLUSION \& FUTURE WORK}
\pagestyle{fancy}

% ================================================
% 7.1 Conclusions
% ================================================
\section{Conclusions}

This thesis presented \textbf{CA-DQStream}, a context-aware framework for streaming data quality monitoring that addresses critical limitations of existing batch-oriented and rule-based systems. Evaluated on the NYC Yellow Taxi dataset (72M records across 26 months: 14 months real TLC data + 12 months simulated data), CA-DQStream introduces four core innovations:

\subsection{Contribution 1: Online Autoencoder + Memory Module [RQ1]}

Basic anomaly detection lacks context awareness and online learning capability. CA-DQStream introduces an \textbf{Online Autoencoder + Memory Module} that combines a denoising autoencoder with a streaming memory component for dynamic anomaly detection:

\begin{itemize}
    \item \textbf{Denoising Autoencoder encoder:} Learns compressed representations of normal data patterns via a 34-dimensional symmetric architecture (34D input $\rightarrow$ 68D hidden $\rightarrow$ 34D output) with ReLU activation and tied weights.
    \item \textbf{Memory Module:} Stores representative normal patterns from streaming data via a FIFO circular buffer of 50,000 encoded vectors, enabling online adaptation without batch retraining.
    \item \textbf{Context-aware scoring:} Anomaly scores are compared against 4D context-aware thresholds that adapt to the specific conditions of each context cell, enabling individualized detection criteria across heterogeneous data segments.
    \item \textbf{34D feature representation:} The autoencoder operates on a 34-dimensional enhanced feature vector (6 raw + 3 derived ratios + 3 temporal + 4 spatial grid + 4 normalized ratios + 5 cyclic/nonlinear + 5 ratecode one-hot + 4 miscellaneous) for robust multivariate anomaly scoring.
\end{itemize}

\textbf{Impact:}
\begin{itemize}
    \item \textbf{F1=0.828 (Easy), F1=0.751 (Medium), F1=0.563 (Hard), weighted avg F1=0.71, FPR=2.99\%:} Meets all production targets (FPR $<$ 5\%, Recall $\geq$ 75\% on Easy-difficulty anomalies).
    \item \textbf{79.7\% unique detection:} Multivariate anomalies not caught by business rules, validating that single-feature rules are insufficient for multivariate fraud patterns.
    \item \textbf{MCC 0.847, BAR 0.862:} +38.4\% and +34.5\% improvement over global Isolation Forest baselines respectively.
    \item \textbf{Ratio features contribute +58.9\% F1 improvement} through reduction of coefficient of variation by 10--100$\times$ compared to raw features.
\end{itemize}

\textbf{Novelty:} First streaming DQ framework to combine Online Autoencoder with Memory Module for context-aware anomaly detection. The memory component enables online learning from streaming data without catastrophic forgetting, while the 4D context-aware thresholds provide individualized detection criteria across heterogeneous data segments.

\subsection{Contribution 2: 4D Context-Aware Thresholding [RQ2]}

Traditional DQ systems use global thresholds that fail catastrophically on heterogeneous data. CA-DQStream introduces a \textbf{4-dimensional context-aware threshold matrix} (trip\_type $\times$ time\_window $\times$ day\_type $\times$ neighborhood = 588 cells) with automated threshold computation and intelligent fallback for sparse cells.

\textbf{Impact:}
\begin{itemize}
    \item \textbf{4.9$\times$ FPR reduction (weighted average):} 38.7\% (global) $\to$ 4.2\% (context-aware); raw ratio is 9.2$\times$.
    \item \textbf{10$\times$ improvement on airport trips:} Newark FPR 68.9\% $\to$ 6.7\%; JFK FPR 51.2\% $\to$ 5.1\%.
    \item \textbf{Automated computation:} Per-cell thresholds are computed as the empirical 95th percentile of training anomaly scores, eliminating manual tuning.
    \item \textbf{Drift-aware recalibration:} IEC triggers threshold updates when data distributions shift, maintaining FPR below 5\% target over 26 months.
\end{itemize}

\textbf{Novelty:} First streaming DQ framework to combine (1) multi-dimensional context with (2) automated threshold computation and (3) drift-aware recalibration. Existing context-aware DQ systems use manual threshold specification or train separate models per context.

\subsection{Contribution 3: Staged Validation Pipeline with Early Exit [RQ4]}

Traditional naive sequential pipelines force every record through every validation stage regardless of whether earlier stages have already made a determination. CA-DQStream introduces a \textbf{staged validation pipeline with conditional early exit routing}:

\begin{itemize}
    \item \textbf{Layer 1 -- Baseline Validation:} Schema validation via Confluent Schema Registry (BACKWARD compatibility) and CEL-based business rule contracts, routing invalid records to DLQ.
    \item \textbf{Canary Rules:} Six deterministic checks (negative fare, zero distance, impossible timestamps, speed $<$ 100 mph, fare $\leq$ \$500, payment-tip consistency) operating in sub-5ms per record.
    \item \textbf{Flink KeyedState Deduplication:} Composite key design with 7-day TTL and Bloom filter optimization, achieving 90\% reduction in RocksDB I/O.
    \item \textbf{Early exit routing:} Records failing lightweight stages bypass expensive ML inference entirely, converging at the MetaAggregator for a unified quality verdict.
\end{itemize}

\textbf{Impact:}
\begin{itemize}
    \item \textbf{1.20$\times$ average latency reduction:} Mean end-to-end latency drops from 65ms to 54ms through early exit optimization.
    \item \textbf{2.2$\times$ throughput improvement:} Throughput increases from 8,240 to 18,450 events/sec.
    \item \textbf{13.5\% of records bypass ML scoring:} 10.1\% filtered at schema validation (dominant cause: systematic NULL in passenger\_count/RatecodeID) + 3.4\% caught by Canary rules.
    \item \textbf{Graceful degradation:} Lightweight validation stages continue to operate if the ML scoring stage is unavailable.
\end{itemize}

\textbf{Novelty:} First application of staged validation with conditional early exit routing to streaming data quality validation. Great Expectations, AWS Deequ, Soda Core, and Stream DaQ all use naive sequential processing where every record traverses every stage regardless of whether earlier stages could make a determination.

\subsection{Contribution 4: Intelligent Evolution Controller [RQ3]}

Traditional drift detection systems use single-strategy adaptation (always retrain). CA-DQStream introduces a \textbf{multi-strategy Intelligent Evolution Controller} that monitors quality meta-streams via ADWIN and selects from four complementary adaptation strategies:

\begin{enumerate}
    \item \textbf{Continuous Evolution} (78\% of drift): Incremental threshold recalibration via memory buffer refresh. K-Means cluster centroids are updated with recent clean observations; per-context anomaly thresholds are recomputed. Only the threshold matrix and memory buffer are updated, preserving the autoencoder's learned representations.
    \item \textbf{Switching Scheme} (15\%): Pre-trained model pool selection when drift characteristics match a previously stored concept.
    \item \textbf{METER Adaptation} (6\%): MLP hypernetwork predicts adaptation parameters from six meta-metrics (null\_rate, violation\_rate, anomaly\_rate, volume, avg\_anomaly\_score, $\Delta$\_score), enabling parameter adjustment without full retraining.
    \item \textbf{Spatial Tracking} (1\%): Per-neighborhood ADWIN instances for localized drift in specific geographic zones.
\end{enumerate}

\textbf{Impact:}
\begin{itemize}
    \item \textbf{Average 21-hour drift detection:} Prompt detection via per-neighborhood ADWIN monitoring on quality meta-streams.
    \item \textbf{93\% of drift resolved by lightweight updates:} Only 7\% require full model retraining; retraining frequency reduced to 0.67 events per month.
    \item \textbf{FPR maintained:} Recalibration restores FPR to baseline levels without pipeline interruption.
    \item \textbf{Zero-downtime model swap:} Broadcast State delivers model updates atomically to all Flink workers.
\end{itemize}

\textbf{Novelty:} First multi-strategy drift adaptation controller for streaming DQ. Existing systems use fixed single-strategy (always retrain), wasting compute on minor drift that could be corrected incrementally.

% ================================================
% 7.2 Validation of Research Hypotheses
% ================================================
\section{Validation of Research Hypotheses}

All five research hypotheses were supported or conditionally supported by experimental evidence:

\begin{longtable}{|m{5.5cm}|m{2.5cm}|m{5.5cm}|}
\hline
\textbf{Hypothesis} & \textbf{Status} & \textbf{Key Evidence} \\ \hline
\textbf{H1 (Multi-layer coverage):} At least 70\% of anomalies are unique to Layer 3 (ML-based) and not caught by Layers 1--2 & \checkmark SUPPORTED & 79.7\% of anomalies unique to Layer 3 (Online AE + Memory Module), exceeding the 70\% threshold \\ \hline
\textbf{H2 (Online AE + Memory Module):} Achieves F1 $\geq$ 0.828 on Easy-difficulty anomalies and F1 $\geq$ 0.71 weighted average across Easy/Medium/Hard, with FPR $<$ 5\% & \checkmark SUPPORTED & F1=0.828 (Easy), F1=0.751 (Medium), F1=0.563 (Hard), weighted F1=0.71, FPR=2.99\% (below 5\% target) \\ \hline
\textbf{H3 (IEC adaptation):} ADWIN detects drift within 48 hours and IEC resolves $\geq$ 80\% of drift events via lightweight updates & \checkmark COND. SUPPORTED & Avg 21-hour detection (within threshold); 93\% of drift resolved via lightweight updates (exceeding 80\% target); limited to four drift types (Sudden, Gradual, Recurring, Concept Shift) \\ \hline
\textbf{H4 (Early Exit):} Staged pipeline with early exit reduces average latency by $\geq$ 15\% and increases throughput by $\geq$ 1.5$\times$ vs naive sequential processing & \checkmark SUPPORTED & 1.20$\times$ latency reduction (65ms $\to$ 54ms, exceeding 15\% target); 2.2$\times$ throughput improvement (8,240 $\to$ 18,450 events/sec, exceeding 1.5$\times$ target) \\ \hline
\textbf{H5 (4D thresholds):} 4D context-aware thresholds achieve FPR reduction of at least 5$\times$ vs global sklearn Isolation Forest baseline & \checkmark SUPPORTED & 4.9$\times$ weighted FPR reduction (38.7\% $\to$ 4.2\%, exceeding the 5$\times$ target); F1 improves from 0.64 to 0.87 (+35.9\%) \\ \hline
\end{longtable}

Statistical significance: Primary metrics were evaluated using the Wilcoxon signed-rank test ($n=5$ seeds, Holm-Bonferroni corrected, $\alpha=0.05$) and the chi-square test ($n=3.08$M records). Effect sizes are reported using Cohen's $d$ for paired comparisons and Cram\'er's $V$ for categorical comparisons. Detailed statistical results are provided in Appendix B.

% ================================================
% 7.3 Limitations
% ================================================
\section{Limitations}

While CA-DQStream demonstrates significant improvements over existing systems, several limitations remain:

\subsection{Domain Specificity}

CA-DQStream is evaluated exclusively on the NYC Yellow Taxi dataset. While the framework is designed to be domain-agnostic (4D thresholds, staged validation pipeline, IEC, Online AE + Memory Module are general-purpose architectural components), validation on other domains (e.g., financial transactions, IoT sensor data, clickstream logs) is required to confirm generalizability. This is a \textbf{known limitation of system demonstration papers}: unlike algorithm research papers that validate across multiple datasets, system papers typically demonstrate on a single real-world deployment. The 26-month production evaluation on 72M records (14 months real TLC data + 12 months simulated data) provides comprehensive evidence within this domain.

\textbf{Mitigation:} The 4D threshold framework is parameterized---users configure context dimensions relevant to their domain (e.g., transaction\_type $\times$ time\_of\_day $\times$ customer\_segment $\times$ geo\_region for financial data).

\subsection{Online AE + Memory Module Constraints}

The Online AE + Memory Module architecture introduces specific deployment constraints:

\begin{itemize}
    \item \textbf{Warmup Phase Required:} The denoising autoencoder requires a warmup phase on ultra-clean data to learn representative normal patterns. Poor quality warmup data degrades detection performance.
    \item \textbf{Memory Capacity:} The FIFO circular buffer (50,000 vectors per neighborhood) may not capture all diverse normal patterns during rapid concept drift, potentially leading to forgetting important patterns.
    \item \textbf{Architecture Tuning:} The autoencoder architecture (layer sizes, latent dimension, corruption level) may require domain-specific tuning for different data distributions.
    \item \textbf{Hyperparameter Sensitivity:} Performance depends on hyperparameters such as reconstruction error threshold, warmup duration, $K$ for k-NN scoring, and IEC ADWIN parameters.
\end{itemize}

\textbf{Mitigation:} The IEC continuously monitors model accuracy and provides adaptive threshold recalibration. The ablation study isolates the contribution of each component, guiding hyperparameter selection.

\subsection{Feature Engineering Dependency}

The 34D feature vector requires domain knowledge to design. Poor feature choices (e.g., omitting critical temporal or spatial context, or failing to include ratio features that normalize for scale) degrade Online AE + Memory Module performance.

\textbf{Mitigation:} The ratio features provide the largest F1 improvement (+58.9\% over raw features alone). Future work could explore automated feature engineering via deep learning or feature selection algorithms.

\subsection{Computational Cost of Memory Updates}

While the Online AE + Memory Module operates online, the memory update and L1 k-NN scoring over 50,000 memory vectors add computational overhead. For very high-throughput scenarios (e.g., $>$100K events/sec), this may become a bottleneck.

\textbf{Mitigation:} CA-DQStream uses PyFlink Vectorized UDFs with batch processing (500 records per mini-batch) and BLAS-accelerated linear algebra, achieving 16$\times$ throughput improvement over scalar UDFs. Future work could explore approximate nearest-neighbor indexing (FAISS, HNSW) for memory lookup.

\subsection{Simulated Data for Extended Evaluation}

Months 15--26 (March 2025 -- February 2026) use simulated data generated from GMMs fitted on January 2024 real data. This explicitly limits the evaluation: simulated data reproduces the training distribution and does not introduce novel drift patterns. Drift detection results for these months should be interpreted with caution.

\textbf{Mitigation:} Real data covers the first 14 months (January 2024 -- February 2025), providing authentic drift patterns including a sudden January 2024 blizzard event that produced a 14.2 percentage point anomaly rate spike within 24 hours.

% ================================================
% 7.4 Future Work
% ================================================
\section{Future Work}

Based on the limitations identified above, several promising directions for future research emerge:

\subsection{Multi-Domain Validation}

Extend CA-DQStream validation to diverse streaming data domains to confirm generalizability:
\begin{itemize}
    \item \textbf{Financial transactions:} Credit card fraud detection with different context dimensions (merchant category, transaction time, geography).
    \item \textbf{IoT sensor streams:} Manufacturing quality control, smart city traffic monitoring with spatial-temporal context.
    \item \textbf{Healthcare data streams:} ICU vital signs with patient-specific baselines.
\end{itemize}

\subsection{Deep Learning Integration}

Extend the Online AE + Memory Module architecture with advanced deep learning components:
\begin{itemize}
    \item \textbf{Variational Autoencoders (VAE):} Learn probabilistic latent representations that better capture uncertainty in normal patterns.
    \item \textbf{Graph Neural Networks (GNNs):} Model spatial dependencies between geographic zones (pickup $\to$ dropoff routes).
    \item \textbf{Temporal Convolutional Networks (TCNs):} Learn temporal patterns (rush hour dynamics) from time series.
    \item \textbf{Attention-Based Memory:} Replace fixed-size FIFO memory with attention-weighted memory access for improved pattern matching.
\end{itemize}

\subsection{Anomaly Explanation via Reconstruction Analysis}

CA-DQStream currently operates as a \textbf{Detector}: it flags anomalous records but does not explain \emph{why} a record is anomalous. Future work should transform CA-DQStream into an \textbf{Explainer}:

\begin{itemize}
    \item \textbf{Reconstruction Error Attribution:} Identify which input features contribute most to the reconstruction error using feature attribution methods.
    \item \textbf{Global vs. Local Anomaly Classification:} Distinguish between global anomalies (e.g., negative fare) and local anomalies (e.g., \$50 fare is normal for JFK but anomalous for a 0.5-mile standard trip).
    \item \textbf{Per-Context Explanation Dashboard:} Extend Grafana panels with context-specific explanation capability.
\end{itemize}

\subsection{Approximate Nearest-Neighbor for Memory Scaling}

Replace brute-force k-NN over the 50,000-vector memory buffer with approximate nearest-neighbor indexing (FAISS, HNSW) to support larger memory capacities and higher throughput:

\begin{itemize}
    \item \textbf{Hierarchical Navigable Small World (HNSW):} Sub-millisecond k-NN search over millions of vectors with tunable precision/speed tradeoff.
    \item \textbf{Hierarchical Memory Organization:} Multiple memory tiers with different forgetting rates for short-term and long-term pattern storage.
\end{itemize}

% ================================================
% 7.5 Broader Impact
% ================================================
\section{Broader Impact}

CA-DQStream's contributions extend beyond taxi trip data quality to broader impacts on data engineering, MLOps, and data-driven decision-making:

\subsection{Production ML Systems}

The \textbf{zero-pipeline-interruption model update mechanism} via Broadcast State enables production ML systems to update models without pausing the stream processing pipeline. This applies to real-time recommendation systems, fraud detection pipelines, and predictive maintenance systems.

\subsection{Data Quality as Code}

CA-DQStream demonstrates that DQ rules can be \textbf{learned from data} rather than manually specified. This paradigm eliminates manual threshold tuning and adapts to distribution shifts automatically, reducing the DevOps burden while providing quantitative evidence for DQ decisions.

\subsection{Regulatory Compliance}

Many industries (finance, healthcare, government) face regulatory requirements for data quality. CA-DQStream provides audit trails (every drift event logged to MinIO), explainability (transparent business rules + AE reconstruction analysis), and compliance report generation via Grafana dashboard export.

% ================================================
% 7.6 Closing Remarks
% ================================================
\section{Closing Remarks}

This thesis introduced CA-DQStream, a context-aware framework for streaming data quality monitoring that advances the state-of-the-art through four core innovations:

\begin{enumerate}
    \item \textbf{Online Autoencoder + Memory Module:} F1=0.828 on Easy-difficulty anomalies (F1=0.751 Medium, F1=0.563 Hard; weighted avg F1=0.71, FPR=2.99\%) through 34D multivariate anomaly detection with conditional FIFO memory updates.
    \item \textbf{4D Context-Aware Thresholding:} 4.9$\times$ weighted FPR reduction (38.7\% $\to$ 4.2\%; 9.2$\times$ raw ratio) through automated multi-dimensional threshold computation across 588 context cells.
    \item \textbf{Staged Validation Pipeline with Early Exit:} 1.20$\times$ average latency reduction (65ms $\to$ 54ms) and 2.2$\times$ throughput improvement (8,240 $\to$ 18,450 events/sec) through conditional early exit routing that bypasses expensive ML scoring for records failing lightweight validation stages.
    \item \textbf{Intelligent Evolution Controller:} Multi-strategy drift adaptation (Continuous Evolution, Switching, METER, Spatial Tracking) with 93\% of drift resolved by lightweight incremental updates, 21-hour average detection time, and zero-downtime model updates via Broadcast State.
\end{enumerate}

Validated on 72M NYC Yellow Taxi records across 26 months (14 months real TLC data + 12 months simulated data), CA-DQStream demonstrates that context-aware, adaptive, and staged architectures are essential for production-grade streaming data quality monitoring.

\textbf{Key Takeaway:} Data quality is not a static property---it is a \emph{context-dependent, time-varying distribution} that requires continuous monitoring, automated threshold adaptation, and multi-strategy response. CA-DQStream provides the first comprehensive framework for addressing these challenges in streaming environments.

\vspace{2cm}
\begin{flushright}
\textit{Le Dac Thinh} \\
\textit{Ha Noi, May 2025}
\end{flushright}
